\section{Two defences that need no training}
\label{sec:defense}

If the mechanism of an evasion is understood, a cheap fix should blunt it. We
test two defences that can be applied after the model is trained, requiring no
fine-tuning, no extra data and no change to the deployed weights, and we measure
them on the \emph{same} stored perturbations the attacks produced, so the
comparison is like for like.

\paragraph{D1: input normalisation.} The text is NFKC-normalised (a standard
Unicode cleanup), invisible and format characters are removed, Cyrillic and Greek
look-alike characters are folded back to their ASCII equivalents, and whitespace
is collapsed. This targets exactly the mechanism of the character attack: the
model sees a different token where a human sees the same word. It should
therefore rescue character-level evasions and do almost nothing for word-level
ones --- and that null result is evidence that the two attack families fail for
different reasons. The transform is idempotent (running it twice changes nothing
the second time), which is checked in \texttt{tests/check\_defense.py}; a
transform that was not idempotent would make the ``defended'' number an artefact
of running it twice.

\paragraph{D2: re-thresholding under a false-positive budget.} The model's
$P(\text{malicious})$ is compared against a threshold $\tau$ instead of the
default $0.5$. $\tau$ is fitted on the clean benign half of the evaluation split
to a 5\% false-positive budget, and it is the \emph{smallest} threshold that meets
that budget, that is, the least conservative setting the budget allows. Two
figures travel with every headline number: the clean false-positive rate at
$\tau$, and the clean true-positive rate it costs. Without the second, a drop in
ASR cannot be told apart from a filter that has simply stopped flagging mail.

The direction of the search matters and is easy to get backwards. A larger
$\tau$ always has a lower false-positive rate, so maximising it under an upper
bound is trivially satisfied by ``above every clean score'' --- a filter that
rejects everything. Our first implementation did exactly that, and because it
also counted a flip whenever an edited text fell below $\tau$, without asking
whether the original was above it, it reported a 100\% ASR for the defence. Both
rules are now enforced in code: $\tau$ is the minimum feasible value, and a sample
counts as flipped only if the \emph{same} rule classified the \emph{original}
text as malicious.

\paragraph{Reading the outcome.} Three outcomes are possible and all three are
informative: a defence that removes the flips it targets (D1 against the
character attack), a defence that removes none (D1 against a word attack it
cannot see), and a defence that trades a small amount of clean precision for a
large reduction in evasions (D2). What we do not do is report a defence at a
different operating point than the baseline and call the difference robustness.

\IfFileExists{tables/t6_defense.tex}{% Generated by src/make_paper_tex.py from results/defense_summary.csv.
Under the char attack, normalisation moves the aggregate ASR from 1.0\% to 0.0\%.
Under the word attack, normalisation moves the aggregate ASR from 2.5\% to 2.0\%.
Under the saliency attack, normalisation moves the aggregate ASR from 7.5\% to 6.9\%.
The threshold is not free: at the 5\% clean false-positive budget its tightest operating point (SmolLM2-1.7B-Instruct, SMS smishing) keeps clean recall at 23.0\%, down from 100.0\% at the default threshold.
1 of 6 model--task cells cannot meet the budget at all: holding their false positives down requires a threshold that rejects every message.
}{}

Table~\ref{tab:defense} is read column by column, because its baseline column
uses the probability channel the threshold needs and the attack tables do not;
the next section says what that costs the comparison.

\begin{table*}[t]
\centering
\small
\caption{Training-free defences, measured on the same stored perturbations. \emph{None} is the undefended ASR, \emph{norm} applies input normalisation only, \emph{thr} replaces the default decision threshold with one fitted on clean benign mail to a 5\% false-positive budget. Rows are pooled over the three models; $n$ is the number of originals the undefended rule itself flagged, and \emph{n/a} means a defence flagged none of them, so there was nothing left to flip. $\dagger$ marks a row containing a model-task pair whose false-positive budget could not be met: its threshold is a rejection rule, not a defence, and the number is there to be ignored.}
\label{tab:defense}
\begin{tabular}{llcccc}
\toprule
\textbf{Mode} & \textbf{Task} & \textbf{None} & \textbf{Norm.} & \textbf{Thr.} & \textbf{$n$} \\
\midrule
 char & SMS smishing & 0.6\% & 0.0\% & 2.3\%$\dagger$ & 515 \\
 char & Phishing email & 1.4\% & 0.0\% & 7.7\% & 486 \\
\midrule
 word & SMS smishing & 2.1\% & 2.1\% & 4.7\%$\dagger$ & 512 \\
 word & Phishing email & 2.9\% & 1.9\% & 5.4\% & 483 \\
\midrule
 saliency & SMS smishing & 6.6\% & 6.6\% & 19.9\%$\dagger$ & 515 \\
 saliency & Phishing email & 8.4\% & 7.3\% & 18.9\% & 486 \\
\bottomrule
\end{tabular}
\end{table*}

